\section*{第 \ref{ch:b1:atmospheric-circulation} 章——大气环流}
% values: calc/ch28.py

\begin{solution}{exo:b1:atmospheric-circulation:1}
(a)~$f = 2\Omega\sin\varphi$：赤道处为 $0$，在 \ang{30} 处为 \qty{7.29e-5}{s^{-1}}，
在 \ang{45} 处为 \qty{1.031e-4}{s^{-1}}，极点处为 \qty{1.458e-4}{s^{-1}}。
(b)~$fV = \num{1.031e-4} \times 10 = \qty{1.03e-3}{m.s^{-2}}$，约为
$g$ 的 \num{1.05e-4} 倍。(c)~北半球的偏转位于运动方向的右侧：
向北吹的风被推向东侧。
\end{solution}

\begin{solution}{exo:b1:atmospheric-circulation:2}
(a)~$|\nabla p|/\rho = (400/\num{4e5})/1.225 = \qty{8.16e-4}{m.s^{-2}}$。
(b)~$10/\num{8.16e-4} = \qty{12250}{s}$，约 \qty{3.4}{h}。(c)~北纬 \ang{45} 处风速
为 \qty{10}{m/s} 的风，其科里奥利加速度为 \qty{1.03e-3}{m.s^{-2}}，二者大小相当
（比值为 $0.79$）。开始运动后数小时之内，空气即受到能与\omterm{def:b1:atmospheric-circulation:pgf}{气压梯度力}平衡的
科里奥利力，风速趋于接近\omterm{def:b1:atmospheric-circulation:geostrophic}{地转风}速，
此处为 $\num{8.16e-4}/\num{1.031e-4} = \qty{7.9}{m/s}$，
而非沿梯度方向持续加速。
\end{solution}

\begin{solution}{exo:b1:atmospheric-circulation:3}
(a)~$f = 2 \times \num{7.292e-5} \times \sin\ang{40} = \qty{9.374e-5}{s^{-1}}$；
$|\vec v_g| = (400/\num{2.5e5})/(1.225 \times \num{9.374e-5}) = \qty{13.9}{m/s}$，或
$13.93/0.5144 = \qty{27.1}{kn}$。(b)~北半球的风使低压位于其左侧：
低压位于南侧时，风向西吹，即东风
（\omterm{def:b1:spherical-astronomy:horizontal}{方位角} \ang{90}）。(c)~风杆自测站圆圈指向东方，朝向
来风方向；\qty{27.1}{kn} 取整为 \qty{25}{kn}，即两根完整羽毛与一根半羽毛，
画在风杆顺时针一侧，此处即南侧，朝向低压。
\end{solution}

\begin{solution}{exo:b1:atmospheric-circulation:4}
(a)~北纬 \ang{15}：东北\omterm{def:b1:atmospheric-circulation:winds}{信风}；南纬 \ang{15}：东南\omterm{def:b1:atmospheric-circulation:winds}{信风}；
北纬 \ang{45} 与南纬 \ang{45}：\omterm{def:b1:atmospheric-circulation:winds}{西风带}。(b)~在 \ang{15} 处，\omterm{def:b1:atmospheric-circulation:cells}{哈得来环流圈}的地表
空气流回赤道，受科里奥利力作用向西偏转
（北半球向右，南半球向左）：为东风，$u$ 为负。在 \ang{45} 处，
\omterm{def:b1:atmospheric-circulation:cells}{费雷尔环流圈}的地表空气流向极地，并向东偏转：为西风，
$u$ 为正。(c)~三个\omterm{def:b1:spherical-astronomy:geographic}{纬度}中，北纬 \ang{30} 处气压最高，
为 \qty{1016.8}{hPa}：\omterm{def:b1:atmospheric-circulation:cells}{副热带高压}位于北纬 \ang{30} 附近（完整剖面显示其极大值位于
北纬 \ang{35} 附近）。
\end{solution}

\begin{solution}{exo:b1:atmospheric-circulation:5}
(a)~$|\nabla p|/\rho = (200/\num{5e5})/1.225 = \qty{3.27e-4}{m.s^{-2}}$，除以
$f$ 得：\ang{60} 处为 \qty{2.6}{m/s}，\ang{50} 处为 \qty{2.9}{m/s}，
\ang{30} 处为 \qty{4.5}{m/s}，\ang{20} 处为 \qty{6.5}{m/s}，\ang{10} 处为 \qty{12.9}{m/s}，
\ang{5} 处为 \qty{25.7}{m/s}。(b)~科里奥利力使风发生偏转所需的
时间量级为 $1/f$：\ang{60} 处约 \qty{2.2}{h}，\ang{10} 处约 \qty{11}{h}，\ang{5} 处约 \qty{22}{h}，
赤道处则无限长。靠近赤道时，气压场会发生变化，摩擦也会起作用，
这远早于风能沿等压线转向之前：平衡
永远无法建立，空气以适中的速度部分地沿梯度方向流动。
\end{solution}

\begin{solution}{exo:b1:atmospheric-circulation:6}
(a)~$u = 464.6\,\sin^2\varphi/\cos\varphi$：\ang{10} 处为 \qty{14.2}{m/s}，
\ang{15} 处为 \qty{32.2}{m/s}，\ang{20} 处为 \qty{57.8}{m/s}，
\ang{30} 处为 \qty{134.1}{m/s}。(b)~$u = 464.6\,(\cos^2\ang{10} - \cos^2\varphi)/\cos\varphi$：
\ang{15} 处为 \qty{17.7}{m/s}，\ang{20} 处为 \qty{42.9}{m/s}，
\ang{30} 处为 \qty{117.9}{m/s}。(c)~起始于离赤道较远处的空气，起始时离自转轴较近，
角动量余量较少，故在各\omterm{def:b1:spherical-astronomy:geographic}{纬度}处西风均较弱。
上升支的扩展使预测的\omterm{def:b1:atmospheric-circulation:jet}{急流}降低，但远不足以
达到观测值 \qty{40}{m/s}。
\end{solution}

\begin{solution}{exo:b1:atmospheric-circulation:7}
(a)~风向东北吹。背对风、面向东北时，
低压位于左侧，即西北方向。(b)~风向
东南吹；在南半球，背对风的人，低压位于其右侧，
即西南方向。(c)~地面风由等压线方向偏向
低压，这一偏转相对于北半球的\omterm{def:b1:atmospheric-circulation:geostrophic}{地转风}为逆时针方向。
因此，\omterm{def:b1:atmospheric-circulation:geostrophic}{地转风}及与其左侧成直角的低压，
相对于(a)的答案顺时针转动：低压位于西北与北之间，
\omterm{def:b1:spherical-astronomy:horizontal}{方位角}为 $\ang{315} + \alpha$，交角为
$\alpha$。
\end{solution}

\begin{solution}{exo:b1:atmospheric-circulation:8}
(a)~赤道与北纬 \ang{30} 之间的面积为 $2\pi a_\oplus^2\sin\ang{30} =
\qty{1.275e14}{m^2}$；两层之间每平方米上方的质量为
$\Delta p/g = \num{9e4}/9.81 = \qty{9174}{kg}$；总质量为
\qty{1.170e18}{kg}。(b)~$\num{1.170e18}/\num{8.7e10} = \qty{1.34e7}{s}$，约
\num{156} 天，即五个月。(c)~$\num{1.170e18}/\num{21.4e10} = \qty{5.47e6}{s}$，
约 \num{63} 天，即两个月。
\end{solution}

\begin{solution}{exo:b1:atmospheric-circulation:9}
(a)~$\num{0.01} \times \num{86400} = \qty{864}{m}$ 每天。(b)~$0.864 \times 9.8 =
\qty{8.5}{K}$ 每天。(c)~空气通过发射红外辐射而降温：
大气因辐射失去的能量多于吸收的能量，空气下沉
时受压缩而增温，补偿了这部分损失。(d)~下沉空气增温，相对饱和而言变得干燥，
云会蒸发，降水无法形成：\omterm{def:b1:atmospheric-circulation:cells}{副热带
高压}是晴朗少雨的地带，陆地上在此形成炎热的沙漠。
\end{solution}

\begin{solution}{exo:b1:atmospheric-circulation:10}
(a)~$V^2 + frV - fr|\vec v_g| = 0$，其中 $fr = \qty{55.85}{m/s}$：
$V = \tfrac12\bigl(-55.85 + \sqrt{55.85^2 + 4 \times 55.85 \times 15}\bigr) =
\qty{12.3}{m/s}$，为\omterm{def:b1:atmospheric-circulation:geostrophic}{地转风}的 \qty{82}{\percent}。(b)~$V^2 - frV +
fr|\vec v_g| = 0$ 仅当 $(fr)^2 \geq 4fr|\vec v_g|$ 时有实根，即
$|\vec v_g| \leq fr/4 = \qty{14.0}{m/s}$。(c)~在高压中心附近，$r$ 很小，
此时上限 $fr/4$ 也很小：气压梯度必须很弱，等压线
相隔较远，风力较小。低压则没有这样的限制：其等压线可在中心附近密集排列，
风力可以很强。
\end{solution}

\begin{solution}{exo:b1:atmospheric-circulation:11}
(a)~当 $\varphi$ 较小时，$\sin\varphi \approx \varphi$，$\cos\varphi \approx 1$，
故 $u \approx \Omega a_\oplus\varphi^2$。(b)~$\ang{10} = \qty{0.1745}{rad}$：
$464.6 \times 0.1745^2 = \qty{14.15}{m/s}$，与精确值 \qty{14.23}{m/s} 相比，
偏小 \qty{0.5}{\percent}。(c)~数值求解 $464.6\sin^2\varphi/\cos\varphi = 27.0$
得 $\varphi = \ang{13.7}$。(d)~该环流圈在不到 \ang{14} 处即与实际
年平均\omterm{def:b1:atmospheric-circulation:jet}{急流}相符：角动量守恒只在深热带描述\omterm{def:b1:atmospheric-circulation:cells}{哈得来环流圈}的
上层分支；向极方向约
\ang{15} 以外，涡旋接管了这一过程。
\end{solution}

\begin{solution}{exo:b1:atmospheric-circulation:12}
(a)~平均为北纬 $31.7/12 = \ang{2.6}$；\omterm{def:b1:spherical-astronomy:geographic}{纬度}范围从南纬 \ang{12.8} 至北纬 \ang{10.9}，即
\ang{23.7}。(b)~最北位置出现在九月，比夏至晚三个月；
最南位置出现在二月，比冬至晚两个月。
这一滞后主要来自热量储存，尤其是海洋
\omterm{def:b1:fluid-envelopes-structure:ocean-layers}{混合层}（见\cref{ch:b1:fluid-envelopes-structure}），它增温和降温都较慢。(c)~纬向
平均是对所有\omterm{def:b1:spherical-astronomy:geographic}{经度}取平均：若在大陆上，雨带大幅摆入
夏季半球，则所有\omterm{def:b1:spherical-astronomy:geographic}{经度}的平均仍可越过赤道，
而某一大洋上的雨带并未越过。
\end{solution}

\begin{solution}{pb:b1:atmospheric-circulation:1}
\begin{enumerate}
  \item $\rho = p/(R_dT) = \num{100600}/(287.05 \times 283.15) =
        \qty{1.238}{kg.m^{-3}}$.
  \item $|\nabla p| = \qty{400}{Pa}/\qty{2.5e5}{m} = \qty{1.6e-3}{Pa.m^{-1}}$.
  \item $|\nabla p|/\rho = \num{1.6e-3}/1.238 = \qty{1.293e-3}{m.s^{-2}}$.
  \item $\Omega = 2\pi/\num{86164.2} = \qty{7.292e-5}{rad.s^{-1}}$;
        $f = 2\Omega\sin\ang{50} = \qty{1.117e-4}{s^{-1}}$.
  \item $|\vec v_g| = \num{1.293e-3}/\num{1.117e-4} = \qty{11.57}{m/s}$.
  \item $11.57/0.514444 = \qty{22.5}{kn}$; $11.57 \times 3.6 = \qty{41.7}{km/h}$.
  \item 时间为 $|\vec v_g|/(|\nabla p|/\rho)$，且 $|\vec v_g| = (|\nabla
        p|/\rho)/f$，故 $1/f = \qty{8951}{s}$，即 \qty{2.49}{h}；又 $2\pi/f =
        \qty{15.6}{h}$。持续一天或更久的天气形势因此
        接近平衡。
  \item \omterm{def:b1:atmospheric-circulation:pgf}{气压梯度力}指向北方，朝向低压。\omterm{def:b1:atmospheric-circulation:geostrophic}{地转
        风}沿等压线吹，低压位于其左侧：向东，
        来自西方，\omterm{def:b1:spherical-astronomy:horizontal}{方位角} \ang{270}：西风。
  \item 背对西风、面向东方时，左侧为北，低压即位于此处：
        该规律得到满足。
  \item 风杆自测站圆圈向西伸出，指向风的来
        向；\qty{22.5}{kn} 取整为 \qty{20}{kn}，即两根完整羽毛，它们位于
        风杆顺时针一侧，即北侧，朝向低压。
  \item 南纬 \ang{50}，$f < 0$：风使低压位于其右侧，故风向
        向西吹，来自东方。
  \item $f(\ang{30}) = \qty{7.292e-5}{s^{-1}}$: \qty{17.7}{m/s}; $f(\ang{60}) =
        \qty{1.263e-4}{s^{-1}}$: \qty{10.2}{m/s}.
  \item $f(\ang{5}) = \qty{1.271e-5}{s^{-1}}$ 给出 \qty{101.7}{m/s}，这样的风从未
        在如此大的梯度下被观测到。转向时间 $1/f$ 为 \qty{22}{h}：
        平衡来不及建立，公式不再适用。
  \item $9.0/11.57 = 0.778$；沿等压线 $9.0\cos\ang{25} = \qty{8.16}{m/s}$，
        垂直等压线指向低压 $9.0\sin\ang{25} = \qty{3.80}{m/s}$。
  \item \omterm{def:b1:atmospheric-circulation:geostrophic}{地转风}吹向\omterm{def:b1:spherical-astronomy:horizontal}{方位角} \ang{90}；向北侧低压偏转 \ang{25}，
        则地面风吹向 \ang{65}，来自
        \ang{245}，即西南偏西。
  \item $\tan\alpha = 1/(f\tau)$ 给出 $\tau = 1/(\num{1.117e-4}\tan\ang{25}) =
        \qty{1.92e4}{s}$，即 \qty{5.3}{h}；模型预测 $11.57\cos\ang{25} =
        \qty{10.5}{m/s}$，大于观测的 \qty{9.0}{m/s}：实际摩擦力
        并非严格与风速成线性关系，在此角度下其减慢
        作用强于模型。
  \item 取半径为 $R = \qty{500}{km}$、高为 $h =
        \qty{1}{km}$: $\dot m = 2\pi R h\rho u_r = 2\pi \times \num{5e5} \times 1000
        \times 1.238 \times 3.80 = \qty{1.48e10}{kg/s}$.
  \item $w = \dot m/(\rho\pi R^2) = 2hu_r/R = 2 \times 1000 \times 3.80/\num{5e5} =
        \qty{0.0152}{m/s}$，约每天 \qty{1310}{m}。
  \item $1.31 \times 9.8 = \qty{12.9}{K}$ 每天：空气很快达到饱和，
        低压区便充满云和雨。在\omterm{def:b1:atmospheric-circulation:cyclone}{反气旋}中，
        近地面空气向外流动，上方空气下沉增温，天空
        转晴。
  \item $\Omega a_\oplus = \num{7.292e-5} \times \num{6.371e6} = \qty{464.6}{m/s}$
        在赤道处，北纬 \ang{30} 处为 $464.6\cos\ang{30} = \qty{402.3}{m/s}$。
  \item $u = 464.6\,\sin^2\varphi/\cos\varphi$：北纬 \ang{20} 处为 \qty{57.8}{m/s}，
        北纬 \ang{30} 处为 \qty{134.1}{m/s}。
  \item $u = 464.6\,(\cos^2\ang{10} - \cos^2\ang{30})/\cos\ang{30} =
        \qty{117.9}{m/s}$.
  \item $134.1/42.3 = 3.17$。中\omterm{def:b1:spherical-astronomy:geographic}{纬度}涡旋，即低压与
        \omterm{def:b1:atmospheric-circulation:cyclone}{反气旋}，从上分支夺取角动量并输往别处；
        环流圈的空气沿一段\omterm{def:b1:spherical-astronomy:geographic}{纬度}带上升，而不是
        在赤道处静止上升。纬向平均还将蜿蜒的\omterm{def:b1:atmospheric-circulation:jet}{急流}
        对所有\omterm{def:b1:spherical-astronomy:geographic}{经度}取平均。
  \item $M/M_0 = (402.3 + 42.3)\cos\ang{30}/464.6 = 0.829$：空气损失了
        \qty{17}{\percent} 的角动量，即其作为自由环
        本可保有的量。
  \item P 处\omterm{def:b1:atmospheric-circulation:geostrophic}{地转风}为 \qty{11.6}{m/s}，来自西方；
        地面实测风为 \qty{9.0}{m/s}，约为其五分之四，并因摩擦
        作用向低压偏转 \ang{25}。
\end{enumerate}
\end{solution}
